\subsection*{A people with no stronger plea than God's fidelity}

The formulary opens after catastrophe. Psalm~73 supplies no inventory of
successes and no timetable for relief. It asks why the flock has been cast
off, calls attention to the poor, and places the case inside a single
imperative: look upon \emph{your} covenant. The Gradual repeats that appeal
after Paul has spoken of the promise to Abraham. This is the Sunday's first
governing relation: divine fidelity precedes the worshipper's achievement and
remains the ground of hope when achievement cannot carry the case.

The reception tradition prevents the word ``covenant'' from becoming vague.
Augustine and Cassiodorus hear the new covenant promised in Jeremiah~31;
Bellarmine hears the old covenant with the fathers and the land; Rupert, as
reported by Guéranger, hears Abraham's promise. Their disagreement need not be
eliminated. Galatians itself holds historical moments together without making
them interchangeable. The promise is made to Abraham and to the one seed,
Christ. A law given 430 years later cannot annul it. The one God stands behind
promise, law, conviction, and fulfillment; the law is not opposed to the
promise, although it cannot itself give life.

Cassiodorus identifies the prayer's spiritual posture when he says the
psalmist remembers the promise because merits fail. This does not deny good
works. The Collect immediately asks for actual growth in faith, hope, and
charity and for love of the divine commands as the way to what is promised.
Trent receives the prayer's virtue-increase petition while teaching increased
justification. The same decree safeguards the genuine merit of works done in
grace precisely because Christ's grace originates and accompanies them.
Promise and obedient growth are therefore neither competitors nor a bargain:
the gift makes the answer possible.

\subsection*{One God remains refuge within changing time}

Paul's ``hand of a mediator'' generated different patristic identifications.
Theodoret reads Moses; Jerome, Chrysostom, Ambrosiaster, Victorinus, and
Augustine read Christ in distinct senses; Aquinas names Christ and retains
Moses as his figure. The matrix is valuable because Paul's focus is not a
detached naming exercise. The law's mediated place is bounded by the unity of
God and ordered toward the promised seed. Salvation is not held in the
sinner's own hand.

The Alleluia and Offertory give that doctrine a prayerful form. God has been
refuge across generations; the worshipper's times rest in God's hands.
Augustine sees Moses as minister of the old covenant and prophet of the new,
then contrasts God's eternity with temporal change. Cassiodorus's ancient
psalter text is the Offertory's own \latin{témpora mea}; he reads the human
times of Christ's life and death as resting in the Father's power. Jerome's
Hebrew Psalter, the medieval Gloss, and Bellarmine confirm that ``times'' is a
serious ancient reading, not a modern paraphrase.

This entrusted temporality is not fatalism. The lepers will act under a
command; the people will offer gifts; communicants will ask to advance.
Placing times in God's hands means that contingency, persecution, sickness,
and death do not hold final sovereignty. It does not make human injustice into
God's command or relieve Christians of practical responsibility.

\subsection*{Faith walks before it sees}

Luke's ten sufferers remain at a distance, raise their voice for mercy, and
receive an instruction rather than an immediate declaration of cure. They go,
and they are cleansed while going. The sequence makes obedience real without
turning it into payment. The command does not disclose a technique; Jesus'
free power meets persons who act on his word before seeing its outcome.

The Fathers and Doctors receive the road in complementary ways. Tertullian
calls the journey obedient faith. Cyril says that Jesus sends healing with
those he sends to the priests. Augustine and Bede read the priestly showing
through the Church's ministry. Bruno, Albert, Bonaventure, and the penitential
tradition distinguish divine cleansing from priestly discernment and use the
healing on the road to hold interior contrition with completed confession.
Those are genuine Catholic receptions, but the literal story remains an
encounter within Israel's law concerning priestly certification of cleansing.

The Collect's request, ``make us love what thou commandest,'' now acquires
concrete shape. Love is not only affection for an abstract norm; it moves on a
word while the end is unseen. Yet the Gospel disallows a cruel inference.
Physical healing is never evidence that one person believed more than another,
and delayed or absent healing never proves a failure of faith. All ten in the
story obey and all ten are cleansed. The decisive difference comes afterward.

\subsection*{Gift seeks return without ceasing to be gift}

One sees, turns back, glorifies God loudly, falls at Jesus' feet, and gives
thanks. He does not cause the cleansing retroactively. His return discloses
the relational end toward which gift points. Augustine warns against claiming
amendment as one's own; Bede distinguishes the cleansing received by ten from
the saving word spoken to faith in grateful operation. Recent papal reception
makes the same distinction without coercion: Christ heals freely, while the
one recognizes that salvation is more than an isolated benefit.

The Secret and Communion extend this return at the altar. The Secret's double
\latin{propitiáre} places both the people and their gifts under divine mercy.
The offering is not a repayment that changes a reluctant God. It is a response
whose offerers still ask pardon. Wisdom then speaks of the manna once given to
Israel; the Communion antiphon turns to God with \latin{dedísti nobis,
Dómine}. The change from ``them'' to ``us'' becomes Eucharistic thanksgiving,
not the erasure of Israel's gift.

Rabanus reads manna as Christ, the bread of life; Bonaventure develops its
every-taste language as sacramental abundance; John Paul~II joins the Wisdom
verse to the Eucharist as pilgrim food and messianic foretaste. Gregory the
Great quotes the same Wisdom clause for a different purpose: divine speech
adapts to hearers as the manna did to varied tastes. That non-Eucharistic
reading supplies a devotional limit. ``Every taste'' describes divine
condescension and abundance, not a promise of felt consolation.

\subsection*{The outsider is not a weapon against Israel}

Luke names the grateful person a Samaritan and a stranger; he does not name
the ethnicity of the nine. Some patristic and later witnesses nevertheless
construct a contrast between one grateful Gentile and nine ungrateful Jews.
That reception must be reported because it has shaped Christian preaching,
but it cannot govern the synthesis. \emph{Nostra aetate}~4 excludes the claim
that Jews are rejected or accursed, and Catholic canonical reading remembers
that the covenant, the law, the psalms, the manna, the priestly instruction,
and Jesus himself all arise within Israel.

Theophylact gives the more fruitful moral sense: foreign origin is no obstacle
to pleasing God, and birth among the holy gives no right to pride. The
Samaritan therefore exposes the wideness of mercy and the poverty of every
self-assured insider. The same discipline governs leprosy typology. Augustine,
Gregory, Bede, and medieval preachers use varied skin as an image for varied
error or sin; such reception may call the Christian to self-examination, but
it is not medicine and must never stigmatize sick or disabled persons.

\subsection*{Growth begins again after reception}

The last prayer repeats the Collect's \latin{augméntum}. At the opening,
faith, hope, and charity are to increase; after Communion, the recipients
themselves ask to advance toward the increase of eternal redemption. Nothing
in the wording implies that Christ's redemptive work or the sacrament is
deficient. The incompleteness belongs to the pilgrim's participation and
fruitfulness.

This makes the Samaritan's final movement decisive. Jesus raises and sends
the man who has returned. Thanksgiving does not freeze him at the place of a
past miracle. Communion likewise nourishes a journey. Promise precedes it,
command shapes it, mercy heals it, gratitude turns it toward the giver, and
the heavenly bread sustains it until faith and hope reach their promised end
in charity.
